\documentclass[10pt,aspectratio=169]{beamer}
\usetheme[progressbar=frametitle]{metropolis}
\usepackage[brazil]{babel}
\usepackage[T1]{fontenc}
\usepackage[utf8]{inputenc}
\usepackage{booktabs,amsmath,xcolor,tikz}
\usetikzlibrary{arrows.meta,positioning}
\definecolor{cforte}{HTML}{1F5A96}
\definecolor{pyforte}{HTML}{C0571B}
\setbeamercolor{frametitle}{bg=cforte,fg=white}
\setbeamercolor{progress bar}{fg=pyforte}
\setbeamercolor{alerted text}{fg=pyforte}
\IfFileExists{materiais_p3/dados.tex}{\newcommand{\TabelaDestino}{%
\begin{tabular}{lrrrr}
\toprule
Grafo & $10$ & $20$ & $30$ & Média (s) \\
\midrule
W1 & $-\infty$ & $-\infty$ & $-\infty$ & 0,354 \\
W2 & 10,47 & 8,48 & 8,71 & 1,227 \\
W3 & 4,63 & 5,42 & 4,67 & 2,992 \\
W4 & 9,04 & 8,83 & 9,23 & 13,721 \\
W5 & 37,85 & 42,63 & 37,86 & 14,821 \\
\bottomrule\end{tabular}
}
\newcommand{\TabelaCalibracao}{%
\begin{tabular}{lrr}
\toprule
Vértices & Bellman--Ford & Floyd--Warshall \\
\midrule
80 & 0,026--0,036 & 0,019--0,020 \\
160 & 0,077--0,124 & 0,065--0,073 \\
320 & 0,182--0,242 & 0,291--0,303 \\
\bottomrule\end{tabular}
}
\newcommand{\MatrizUm}{%
\begin{tabular}{lrrrrr}
\toprule
 & 10 & 20 & 30 & 40 & 50 \\
\midrule
10 & $-\infty$ & $-\infty$ & $-\infty$ & $-\infty$ & $-\infty$ \\
20 & $-\infty$ & $-\infty$ & $-\infty$ & $-\infty$ & $-\infty$ \\
30 & $-\infty$ & $-\infty$ & $-\infty$ & $-\infty$ & $-\infty$ \\
40 & $-\infty$ & $-\infty$ & $-\infty$ & $-\infty$ & $-\infty$ \\
50 & $-\infty$ & $-\infty$ & $-\infty$ & $-\infty$ & $-\infty$ \\
\bottomrule\end{tabular}
}
\newcommand{\MatrizDois}{%
\begin{tabular}{lrrrrr}
\toprule
 & 10 & 20 & 30 & 40 & 50 \\
\midrule
10 & 0,00 & 11,73 & 8,14 & 9,49 & 7,08 \\
20 & 3,84 & 0,00 & 7,03 & 8,30 & 4,70 \\
30 & 5,79 & 12,55 & 0,00 & 8,30 & 6,33 \\
40 & 5,05 & 12,32 & 5,36 & 0,00 & 5,71 \\
50 & 4,48 & 12,23 & 7,79 & 8,26 & 0,00 \\
\bottomrule\end{tabular}
}
\newcommand{\MatrizTres}{%
\begin{tabular}{lrrrrr}
\toprule
 & 10 & 20 & 30 & 40 & 50 \\
\midrule
10 & 0,00 & 4,99 & 5,71 & 7,49 & 4,31 \\
20 & 3,48 & 0,00 & 6,19 & 8,27 & 4,32 \\
30 & 3,96 & 5,80 & 0,00 & 8,30 & 3,17 \\
40 & 3,58 & 5,59 & 6,10 & 0,00 & 3,84 \\
50 & 3,25 & 5,41 & 6,12 & 8,10 & 0,00 \\
\bottomrule\end{tabular}
}
\newcommand{\MatrizQuatro}{%
\begin{tabular}{lrrrrr}
\toprule
 & 10 & 20 & 30 & 40 & 50 \\
\midrule
10 & 0,00 & 10,52 & 10,15 & 10,69 & 8,24 \\
20 & 7,79 & 0,00 & 9,19 & 9,67 & 7,58 \\
30 & 7,61 & 10,25 & 0,00 & 8,77 & 6,55 \\
40 & 7,94 & 8,65 & 9,22 & 0,00 & 7,75 \\
50 & 7,99 & 11,41 & 8,80 & 8,55 & 0,00 \\
\bottomrule\end{tabular}
}
\newcommand{\MatrizCinco}{%
\begin{tabular}{lrrrrr}
\toprule
 & 10 & 20 & 30 & 40 & 50 \\
\midrule
10 & 0,00 & 40,20 & 39,56 & 41,55 & 44,91 \\
20 & 39,17 & 0,00 & 39,98 & 49,82 & 45,16 \\
30 & 34,46 & 42,64 & 0,00 & 47,04 & 43,85 \\
40 & 38,73 & 43,91 & 47,30 & 0,00 & 45,24 \\
50 & 44,06 & 40,62 & 45,25 & 49,11 & 0,00 \\
\bottomrule\end{tabular}
}

% Memoria das matrizes densas, obtida de (V+1)^2 celulas.
\newcommand{\TabelaRecursos}{%
\begin{tabular}{lrr}
\toprule
Grafo & Floyd--Warshall & Bellman--Ford denso\\
\midrule
W1, W2 & 7,50 GB & 5,00 GB\\
W3 & 120 GB & 80 GB\\
W4 & 12 TB & 8 TB\\
W5 & 1,2 PB & 0,8 PB\\
\bottomrule
\end{tabular}%
}
}{\newcommand{\TabelaDestino}{%
\begin{tabular}{lrrrr}
\toprule
Grafo & $10$ & $20$ & $30$ & Média (s) \\
\midrule
W1 & $-\infty$ & $-\infty$ & $-\infty$ & 0,354 \\
W2 & 10,47 & 8,48 & 8,71 & 1,227 \\
W3 & 4,63 & 5,42 & 4,67 & 2,992 \\
W4 & 9,04 & 8,83 & 9,23 & 13,721 \\
W5 & 37,85 & 42,63 & 37,86 & 14,821 \\
\bottomrule\end{tabular}
}
\newcommand{\TabelaCalibracao}{%
\begin{tabular}{lrr}
\toprule
Vértices & Bellman--Ford & Floyd--Warshall \\
\midrule
80 & 0,026--0,036 & 0,019--0,020 \\
160 & 0,077--0,124 & 0,065--0,073 \\
320 & 0,182--0,242 & 0,291--0,303 \\
\bottomrule\end{tabular}
}
\newcommand{\MatrizUm}{%
\begin{tabular}{lrrrrr}
\toprule
 & 10 & 20 & 30 & 40 & 50 \\
\midrule
10 & $-\infty$ & $-\infty$ & $-\infty$ & $-\infty$ & $-\infty$ \\
20 & $-\infty$ & $-\infty$ & $-\infty$ & $-\infty$ & $-\infty$ \\
30 & $-\infty$ & $-\infty$ & $-\infty$ & $-\infty$ & $-\infty$ \\
40 & $-\infty$ & $-\infty$ & $-\infty$ & $-\infty$ & $-\infty$ \\
50 & $-\infty$ & $-\infty$ & $-\infty$ & $-\infty$ & $-\infty$ \\
\bottomrule\end{tabular}
}
\newcommand{\MatrizDois}{%
\begin{tabular}{lrrrrr}
\toprule
 & 10 & 20 & 30 & 40 & 50 \\
\midrule
10 & 0,00 & 11,73 & 8,14 & 9,49 & 7,08 \\
20 & 3,84 & 0,00 & 7,03 & 8,30 & 4,70 \\
30 & 5,79 & 12,55 & 0,00 & 8,30 & 6,33 \\
40 & 5,05 & 12,32 & 5,36 & 0,00 & 5,71 \\
50 & 4,48 & 12,23 & 7,79 & 8,26 & 0,00 \\
\bottomrule\end{tabular}
}
\newcommand{\MatrizTres}{%
\begin{tabular}{lrrrrr}
\toprule
 & 10 & 20 & 30 & 40 & 50 \\
\midrule
10 & 0,00 & 4,99 & 5,71 & 7,49 & 4,31 \\
20 & 3,48 & 0,00 & 6,19 & 8,27 & 4,32 \\
30 & 3,96 & 5,80 & 0,00 & 8,30 & 3,17 \\
40 & 3,58 & 5,59 & 6,10 & 0,00 & 3,84 \\
50 & 3,25 & 5,41 & 6,12 & 8,10 & 0,00 \\
\bottomrule\end{tabular}
}
\newcommand{\MatrizQuatro}{%
\begin{tabular}{lrrrrr}
\toprule
 & 10 & 20 & 30 & 40 & 50 \\
\midrule
10 & 0,00 & 10,52 & 10,15 & 10,69 & 8,24 \\
20 & 7,79 & 0,00 & 9,19 & 9,67 & 7,58 \\
30 & 7,61 & 10,25 & 0,00 & 8,77 & 6,55 \\
40 & 7,94 & 8,65 & 9,22 & 0,00 & 7,75 \\
50 & 7,99 & 11,41 & 8,80 & 8,55 & 0,00 \\
\bottomrule\end{tabular}
}
\newcommand{\MatrizCinco}{%
\begin{tabular}{lrrrrr}
\toprule
 & 10 & 20 & 30 & 40 & 50 \\
\midrule
10 & 0,00 & 40,20 & 39,56 & 41,55 & 44,91 \\
20 & 39,17 & 0,00 & 39,98 & 49,82 & 45,16 \\
30 & 34,46 & 42,64 & 0,00 & 47,04 & 43,85 \\
40 & 38,73 & 43,91 & 47,30 & 0,00 & 45,24 \\
50 & 44,06 & 40,62 & 45,25 & 49,11 & 0,00 \\
\bottomrule\end{tabular}
}

% Memoria das matrizes densas, obtida de (V+1)^2 celulas.
\newcommand{\TabelaRecursos}{%
\begin{tabular}{lrr}
\toprule
Grafo & Floyd--Warshall & Bellman--Ford denso\\
\midrule
W1, W2 & 7,50 GB & 5,00 GB\\
W3 & 120 GB & 80 GB\\
W4 & 12 TB & 8 TB\\
W5 & 1,2 PB & 0,8 PB\\
\bottomrule
\end{tabular}%
}
}
\title{Teoria dos Grafos --- COS242}
\subtitle{Parte 3 --- Caminhos mínimos em grafos direcionados}
\author{Gabriel Pires \and Gabriella Lima}
\institute{Universidade Federal do Rio de Janeiro}
\date{2026/2}
\begin{document}
\begin{frame}[plain,noframenumbering] % 0:20
\vfill
{\LARGE\bfseries\color{cforte}\inserttitle\par}
\medskip
{\large\insertsubtitle\par}
\medskip
{\color{pyforte}\rule{\textwidth}{.7pt}\par}
\medskip
Gabriel Pires \quad Gabriella Lima\par
\insertinstitute\par
\insertdate
\vfill
\end{frame}

\begin{frame}{Objetivo e representação} % 0:50
\begin{columns}[T,onlytextwidth]
\column{.56\textwidth}
\begin{itemize}
\item Ler arcos ponderados e preservar o modo não direcionado.
\item Reutilizar BFS, DFS e Dijkstra pela interface de vizinhança.
\item Obter caminhos até um destino e entre todos os pares.
\end{itemize}
\column{.40\textwidth}
\centering
\begin{tikzpicture}[node distance=1.7cm,vert/.style={circle,draw=cforte,thick,minimum size=8mm},>=Stealth]
\node[vert] (u) {\(u\)};
\node[vert,right=of u] (v) {\(v\)};
\draw[->,thick,cforte] (u) -- node[above] {\(w\)} (v);
\end{tikzpicture}\par
\medskip
\small \(u,v,w\) entra na lista de saídas de \(u\).\par
CSR usa deslocamentos e vetores contíguos.
\end{columns}
\end{frame}

\begin{frame}{Bellman--Ford até um destino} % 1:15
\begin{columns}[T,onlytextwidth]
\column{.56\textwidth}
\textbf{Aula 14}
\begin{enumerate}
\item Um vetor de distâncias.
\item Só expande a fronteira atualizada.
\item Para quando a fronteira esvazia.
\end{enumerate}
\column{.41\textwidth}
\centering
\begin{tikzpicture}[>=Stealth,vert/.style={circle,draw=cforte,thick,minimum size=7mm}]
\node[vert] at (0,0) (v) {\(v\)};
\node[vert] at (3,0) (t) {\(t\)};
\draw[->,thick,cforte] (v) to[bend left=28] node[above]{original} (t);
\draw[->,thick,pyforte] (t) to[bend left=28] node[below]{transposto} (v);
\end{tikzpicture}
\end{columns}
\medskip
\[
  d(v)\gets\min\{d(v),\,w(v,u)+d_{\mathrm{anterior}}(u)\}.
\]
\small Uma melhora no passo \(V\) sinaliza ciclo negativo que chega ao destino.
A consulta global usa uma raiz virtual para cobrir todo o grafo.
\end{frame}

\begin{frame}{Floyd--Warshall usa caminhos intermediários} % 1:00
\begin{center}
\begin{tikzpicture}[>=Stealth,vert/.style={circle,draw=cforte,thick,minimum size=7mm}]
\node[vert] at (0,0) (u) {\(u\)};
\node[vert] at (2.5,1) (k) {\(k\)};
\node[vert] at (5,0) (v) {\(v\)};
\draw[->,thick,pyforte] (u) -- node[above left]{\(D_{uk}\)} (k);
\draw[->,thick,pyforte] (k) -- node[above right]{\(D_{kv}\)} (v);
\draw[->,thick,cforte] (u) -- node[below]{\(D_{uv}\)} (v);
\end{tikzpicture}
\end{center}
\[
  D_{uv}\gets\min(D_{uv},D_{uk}+D_{kv}).
\]
\small Matrizes reutilizadas a cada fase. Predecessores reconstroem rotas.
Diagonal negativa indica ciclo. A marcação de \(-\infty\) depende do alcance.
\end{frame}

\begin{frame}{Consulta oficial até o vértice 100} % 0:40
\centering{\large\TabelaDestino}
\par\medskip
\small Bellman--Ford no grafo transposto. Cada tempo é a média de dez rodadas.
Os 12 valores finitos foram conferidos por Dijkstra em C.
\end{frame}

\begin{frame}{Cinco matrizes seletivas de distâncias} % 1:00
\begin{columns}[T,onlytextwidth]
\column{.48\textwidth}\centering \textbf{W1}\par\small\MatrizUm
\column{.48\textwidth}\centering \textbf{W2}\par\small\MatrizDois
\end{columns}
\medskip
\small
\begin{itemize}
\item Bellman--Ford por destino gerou as cinco matrizes oficiais. Floyd--Warshall em escala oficial segue pendente.
\item Em W1, todos os pares podem usar o ciclo \(9931\to12129\to9931\) de peso \(-3{,}88\).
\item Dijkstra em C conferiu as 100 células de W2 a W5.
\end{itemize}
\end{frame}

\begin{frame}{Calibração medida em grafos sintéticos} % 1:05
\centering{\large\TabelaCalibracao}
\par\medskip
\small
\begin{itemize}
\item Dois graus de saída, 4 e 16. Dez rodadas por tamanho, densidade e algoritmo.
\item Faixas das médias em segundos por rodada. Pesos positivos.
\item A posição relativa mudou entre 160 e 320 vértices nesses casos.
\end{itemize}
\end{frame}

\begin{frame}{Todos os pares nos grafos oficiais} % 1:00
\centering\TabelaRecursos
\par\medskip
\begin{itemize}
\item A máquina local tem 4 GB de RAM. As matrizes densas de W1 e W2 já excedem esse total.
\item Projeção para dez rodadas de Floyd em W1 ou W2: ordem de semanas na máquina local.
\item Bellman--Ford em fluxo conserva uma coluna por vez.
\end{itemize}
\end{frame}

\begin{frame}{Conclusões} % 0:50
\begin{columns}[T,onlytextwidth]
\column{.49\textwidth}
\textbf{Contribuições}
\begin{itemize}
\item Grafos direcionados ponderados com reutilização de BFS, DFS e Dijkstra.
\item Bellman--Ford com fronteira ativa e Floyd--Warshall com matrizes em lugar.
\end{itemize}
\column{.47\textwidth}
\textbf{Validação}
\begin{itemize}
\item As 125 distâncias seletivas foram conferidas por Dijkstra em C ou por um certificado de ciclo negativo.
\item A detecção distingue existência global e pares afetados.
\end{itemize}
\end{columns}
\medskip
\alert{Limite experimental.} A comparação de desempenho nos grafos oficiais aguarda dez rodadas de todos os pares por método.
\end{frame}
\end{document}
